\subsection{Eigenvalues by the power method}
\label{sec:eigen}
Every eigenvalue in the pipeline comes from our implementation of the power method as taught in
the course lecture~\cite{tahmid_eigen} (\texttt{core/sirlab/linalg/eigen.py}). Starting from
$x_0 = (1, \dots, 1)^\top$, each iteration forms $y = A x_k$, divides it by its largest-magnitude
entry, and takes that normalizing factor as the eigenvalue estimate. The lecture stops once the
factor's relative change $|\varepsilon_a|$ falls below a tolerance ($10^{-12}$). We additionally
require the normalized vector to stop changing, because the factor alone can settle while $x_k$ is
still far from the eigenvector. On the lecture's $4\times4$ example the largest entry equals 10
from the first step, and deflating that unconverged vector would corrupt every later eigenvalue.
The inverse power method runs the same loop on $A^{-1}$ to find the smallest eigenvalue, solving
$A y = x_k$ with our LU factorization, factored once. Hotelling deflation,
$A_2 = A - \lambda_1 \hat v_1 \hat v_1^\top$, then gives all eigenpairs of a symmetric matrix one at a
time. Our code reproduces the lecture's three worked tables digit for digit
(\texttt{tests/test\_eigen.py}). We use these routines in three places:
\begin{itemize}
  \item the axes of the asymptotic confidence ellipse, from $\mathrm{Cov}(\hat\theta)$ (E11, E14),
    and the cost-valley direction, from $J^\top J$ (E05): power method with Hotelling deflation;
  \item the condition number of $J$ (E10): $\kappa_2(J)^2 = \lambda_{\max}/\lambda_{\min}$ for the
    eigenvalues of $J^\top J$, with $\lambda_{\max}$ from the power method and $\lambda_{\min}$ from
    the inverse power method;
  \item the stability metric $h\,\rho(J_y f)$, where $\rho$ is the largest eigenvalue modulus of
    the SIR Jacobian (E03).
\end{itemize}

\paragraph{When the power method fails.} The method assumes a single dominant eigenvalue,
$|\lambda_1| > |\lambda_2|$, and the SIR Jacobian violates that assumption. Its eigenvalues are 0 and those of the
block $\left[\begin{smallmatrix} -a & -b \\ a & b-\gamma \end{smallmatrix}\right]$, with
$a = \beta I/N$ and $b = \beta S/N$. At the epidemic peak $b = \gamma$, so
\begin{equation}
  \lambda = \tfrac{1}{2}\big(-a \pm \sqrt{a^2 - 4a\gamma}\,\big), \label{eq:jaceig}
\end{equation}
a complex-conjugate pair of equal modulus whenever $a < 4\gamma$. That holds at the peak of every
epidemic with $R_0 < 6.94$, including all three of our regimes ($R_0 = 1.5$, 3, 6). The normalizing
factor then oscillates and never converges. When the iteration stalls like this, the iterate $x$
still settles into the 2-D invariant subspace of the two largest eigenvalues. We orthonormalize
$x$ and $Ax$ (Gram--Schmidt) into $Q$, form the $2\times2$ matrix $H = Q^\top A Q$, and solve its
characteristic equation $\det(H - \lambda I) = 0$ as in the lecture, which gives
$|\lambda| = \sqrt{\det H}$ for a complex pair. We accept the result only once $AQ = QH$ holds to
$10^{-10}$, i.e.\ once the subspace is invariant, and otherwise keep iterating. The same projection
resolves a symmetric matrix whose two largest eigenvalues are $\lambda$ and $-\lambda$. NumPy's
LAPACK-based eigensolvers appear only in the tests, as an independent cross-check.
